\documentclass[12pt, a4paper]{ctexart}
\usepackage{fontspec}
\usepackage{xcolor}
\usepackage{hyperref}
\usepackage{geometry}
\usepackage{enumitem}
\usepackage{parskip}
\usepackage{titlesec}
\usepackage{tcolorbox}
\tcbuselibrary{breakable}
\tcbset{autoparskip}

\usepackage{setspace}
\usepackage{listings}
\usepackage{amssymb}

\geometry{left=2.5cm, right=2.5cm, top=2.5cm, bottom=2.5cm}

\setmainfont{Times New Roman}
\setCJKmainfont{SimSun}

\hypersetup{
    colorlinks=true,
    linkcolor=blue!70!black,
    urlcolor=cyan,
    pdftitle={Java编程基础-逐题精讲解义},
    pdfauthor={专升本-fifine老师}
}

\definecolor{myblue}{RGB}{30,100,200}
\definecolor{lightblue}{RGB}{230,240,255}
\definecolor{darkblue}{RGB}{0,50,120}

\newtcolorbox{bluebox}[1][]{
    breakable,
    colback=lightblue,
    colframe=myblue,
    coltitle=darkblue,
    fonttitle=\bfseries,
    title={#1},
    boxrule=0.8pt,
    arc=4pt,
    left=6pt,
    right=6pt,
    top=6pt,
    bottom=6pt,
    before skip=8pt,
    after skip=8pt
}

\usepackage{etoolbox}
\BeforeBeginEnvironment{bluebox}{\par\vfill}%
\AfterEndEnvironment{bluebox}{\par\vfill}%

\titleformat{\section}
  {\normalfont\Large\bfseries\color{myblue}}{\thesection}{1em}{}
\titlespacing*{\section}{0pt}{3.5ex plus 1ex minus .2ex}{1.5ex plus .2ex}

\title{\textcolor{myblue}{\textbf{Java编程基础 - 逐题精讲解义}}}
\author{专升本-fifine老师 教学组}
\date{2026年03月02日}

\lstset{
    language=Java,
    basicstyle=\ttfamily\small,
    keywordstyle=\color{blue},
    commentstyle=\color{green!60!black},
    stringstyle=\color{orange},
    numbers=left,
    numberstyle=\tiny\color{gray},
    frame=single,
    breaklines=true,
    columns=fullflexible,
    showstringspaces=false
}

\newcommand{\code}[1]{\texttt{#1}}

\begin{document}

\maketitle
\thispagestyle{empty}
\newpage

\section{Java编程基础}

\begin{enumerate}[label=\textbf{\arabic*.}]

	\item \textbf{在Java中，以下哪个关键字用于定义常量？}
	      \begin{enumerate}[label=\Alph*.]
		      \item final
		      \item static
		      \item abstract
		      \item interface
	      \end{enumerate}

	      \begin{bluebox}{答案与深度解析}
		      \textbf{答案：A. final}

		      \textbf{【核心认知框架>}
		      \begin{itemize}[label={}, leftmargin=*, nosep]
			      \item \textbf{题目本质}：考查Java中实现不可变性的关键字
			      \item \textbf{关键区分点}：final保证不可变性，static仅控制作用域
			      \item \textbf{命题人意图}：测试对关键字本质功能的理解
		      \end{itemize}

		      \textbf{【选项原理深度拆解】}

		      \begin{itemize}[leftmargin=*, nosep, label={}]
			      \item[A] \textbf{final — 正确（不可变性）}
			      \begin{itemize}[label={}, leftmargin=*, nosep]
				      \item \textbf{字节码层面}：final字段有写保护机制
				      \item \textbf{JMM保障}：final字段在Java内存模型中提供特殊保障
				      \item \textbf{初始化规则}：实例常量必须在构造器中初始化
				      \item \textbf{代码示例}：
				            \begin{lstlisting}[language=Java]
final int MAX = 100;
// MAX = 200; // 编译错误！
				            \end{lstlisting}
			      \end{itemize}

			      \item[B] \textbf{static — 错误（仅控制作用域）}
			      \begin{itemize}[label={}, leftmargin=*, nosep]
				      \item static仅控制所属权（类级别），与可变性无关
				      \item 类加载初始化线程安全，但static变量可修改
				      \item \textbf{反例}：static int x = 100; x = 200; //编译通过
			      \end{itemize}

			      \item[C] \textbf{abstract — 错误（语法禁止）}
			      \begin{itemize}[label={}, leftmargin=*, nosep]
				      \item abstract禁止修饰变量（JLS §8.3）
				      \item 只能修饰类和方法
			      \end{itemize}

			      \item[D] \textbf{interface — 错误（默认行为）}
			      \begin{itemize}[label={}, leftmargin=*, nosep]
				      \item 接口字段默认为public static final
				      \item 但关键字是final，不是interface
			      \end{itemize}
		      \end{itemize}

		      \textbf{【应背诵知识点>}
		      \begin{itemize}[label={}, nosep]
			      \item final基本类型绝对不可变，引用类型引用地址不可变
			      \item static必须与final组合才能定义常量
			      \item 选项辨析：final=定义常量，static=类级别，abstract=抽象，interface=接口定义
		      \end{itemize}

		      \textbf{【命题趋势与实战策略】}

		      \textbf{考场秒杀}：看到"定义常量"→选final
	      \end{bluebox}

	      \vspace{1em}

	\item \textbf{以下哪个是Java中的合法标识符？}
	      \begin{enumerate}[label=\Alph*.]
		      \item 123abc
		      \item @abc
		      \item \_abc
		      \item class
	      \end{enumerate}

	      \begin{bluebox}{答案与深度解析}
		      \textbf{答案：C. \_abc}

		      \textbf{【核心认知框架>}
		      \begin{itemize}[label={}, leftmargin=*, nosep]
			      \item \textbf{题目本质}：考查Java标识符命名规则
			      \item \textbf{关键区分点}：首字符规则、关键字限制
		      \end{itemize}

		      \textbf{【选项原理深度拆解】}

		      \begin{itemize}[leftmargin=*, nosep, label={}]
			      \item[A] \textbf{123abc — 错误（数字不能作为首字符）}
			      \begin{itemize}[label={}, nosep]
				      \item JLS §3.8：首字符必须是字母、\$\_、Unicode字符
				      \item 数字只能出现在后续位置
			      \end{itemize}

			      \item[B] \textbf{@abc — 错误（@是注解专用字符）}
			      \begin{itemize}[label={}, nosep]
				      \item @是Java保留用于注解的特殊字符
				      \item 不能作为普通标识符使用
			      \end{itemize}

			      \item[C] \textbf{\_abc — 正确（下划线合法）}
			      \begin{itemize}[label={}, nosep]
				      \item 下划线是Java标识符合法首字符
				      \item JDK9+单个\_是保留标识符，但\_abc合法
			      \end{itemize}

			      \item[D] \textbf{class — 错误（关键字禁用）}
			      \begin{itemize}[label={}, nosep]
				      \item class是保留关键字，不能用作标识符
			      \end{itemize}
		      \end{itemize}

		      \textbf{【应背诵知识点>}
		      \begin{itemize}[label={}, nosep]
			      \item 标识符规则：字母、\$\_开头，后续可以是数字、字母、\$\_
			      \item 不能以数字开头
			      \item 不能使用Java关键字
		      \end{itemize}
	      \end{bluebox}

	      \vspace{1em}

	\item \textbf{Java中用于声明整数类型变量的关键字是}
	      \begin{enumerate}[label=\Alph*.]
		      \item int
		      \item float
		      \item double
		      \item long
	      \end{enumerate}

	      \begin{bluebox}{答案与深度解析}
		      \textbf{答案：A. int}

		      \textbf{【核心认知框架>}
		      \begin{itemize}[label={}, leftmargin=*, nosep]
			      \item \textbf{题目本质}：考查Java基本数据类型中整数类型分类
		      \end{itemize}

		      \textbf{【选项原理深度拆解】}

		      \begin{itemize}[leftmargin=*, nosep, label={}]
			      \item[A] \textbf{int — 正确（基本整数类型）}
			      \begin{itemize}[label={}, nosep]
				      \item JVM规范：32位有符号整数
				      \item 默认整数类型
				      \item 取值范围：-2,147,483,648 ~ 2,147,483,647
			      \end{itemize}

			      \item[B] \textbf{float — 错误（浮点类型）}
			      \begin{itemize}[label={}, nosep]
				      \item 32位单精度浮点数
			      \end{itemize}

			      \item[C] \textbf{double — 错误（浮点类型）}
			      \begin{itemize}[label={}, nosep]
				      \item 64位双精度浮点数
			      \end{itemize}

			      \item[D] \textbf{long — 错误（长整型）}
			      \begin{itemize}[label={}, nosep]
				      \item 64位长整型，是整数类型但不是"默认"整数类型
				      \item 使用时需要加L后缀
			      \end{itemize}
		      \end{itemize}

		      \textbf{【应背诵知识点>}
		      \begin{itemize}[label={}, nosep]
			      \item 整数类型：byte(8位)、short(16位)、int(32位)、long(64位)
			      \item 浮点类型：float(32位)、double(64位)
			      \item 默认：整数默认int，浮点默认double
		      \end{itemize}
	      \end{bluebox}

	      \vspace{1em}

	\item \textbf{以下代码的输出结果是}
	      \begin{lstlisting}[language=Java]
int a = 5;
int b = 2;
System.out.println(a / b);
	      \end{lstlisting}
	      \begin{enumerate}[label=\Alph*.]
		      \item 2.5
		      \item 2
		      \item 3
		      \item 2.0
	      \end{enumerate}

	      \begin{bluebox}{答案与深度解析}
		      \textbf{答案：B. 2}

		      \textbf{【核心认知框架>}
		      \begin{itemize}[label={}, leftmargin=*, nosep]
			      \item \textbf{题目本质}：整数除法运算规则——整除取整
		      \end{itemize}

		      \textbf{【选项原理深度拆解】}

		      \begin{itemize}[leftmargin=*, nosep, label={}]
			      \item[A] \textbf{2.5 — 错误（浮点除法结果）}
			      \begin{itemize}[label={}, nosep]
				      \item 需要至少一个操作数是浮点类型
				      \item 改为：a / (double) b 或 a * 1.0 / b
			      \end{itemize}

			      \item[B] \textbf{2 — 正确（向零取整）}
			      \begin{itemize}[label={}, nosep]
				      \item JVM规范：整数除法结果向零取整
				      \item 5 / 2 = 2（余数1被丢弃）
				      \item 负数：-5 / 2 = -2（不是-3）
			      \end{itemize}

			      \item[C] \textbf{3 — 错误（不是四舍五入）}
			      \begin{itemize}[label={}, nosep]
				      \item 四舍五入需要：Math.round(a * 1.0 / b)
			      \end{itemize}

			      \item[D] \textbf{2.0 — 错误（类型不会自动转换）}
			      \begin{itemize}[label={}, nosep]
				      \item int除int结果还是int
				      \item 不会自动转为double
			      \end{itemize}
		      \end{itemize}

		      \textbf{【应背诵知识点>}
		      \begin{itemize}[label={}, nosep]
			      \item 整数相除：截断小数部分（向零取整）
			      \item 取模：5 % 2 = 1
			      \item 浮点除法：至少一个操作数转为浮点类型
		      \end{itemize}
	      \end{bluebox}

	      \vspace{1em}

	\item \textbf{以下哪个是Java中的逻辑与运算符？}
	      \begin{enumerate}[label=\Alph*.]
		      \item \&\&
		      \item ||
		      \item !
		      \item \&
	      \end{enumerate}

	      \begin{bluebox}{答案与深度解析}
		      \textbf{答案：A. \&\&}

		      \textbf{【核心认知框架>}
		      \begin{itemize}[label={}, leftmargin=*, nosep]
			      \item \textbf{题目本质}：区分短路逻辑运算符 vs 按位运算符
		      \end{itemize}

		      \textbf{【选项原理深度拆解】}

		      \begin{itemize}[leftmargin=*, nosep, label={}]
			      \item[A] \textbf{\&\& — 正确（短路逻辑与）}
			      \begin{itemize}[label={}, nosep]
				      \item 第一个操作数为false时，不计算第二个操作数
				      \item 短路特性可提高效率
			      \end{itemize}

			      \item[B] \textbf{|| — 错误（逻辑或）}
			      \begin{itemize}[label={}, nosep]
				      \item 逻辑或运算符
			      \end{itemize}

			      \item[C] \textbf{! — 错误（逻辑非，一元运算符）}
			      \begin{itemize}[label={}, nosep]
				      \item 是一元运算符，不是二元运算符
			      \end{itemize}

			      \item[D] \textbf{\& — 错误（按位与）}
			      \begin{itemize}[label={}, nosep]
				      \item \&是按位与，不是逻辑与
				      \item \&不短路，两边都计算
			      \end{itemize}
		      \end{itemize}

		      \textbf{【应背诵知识点>}
		      \begin{itemize}[label={}, nosep]
			      \item 短路逻辑运算符：\&\&（与）、||（或）
			      \item 按位运算符：\texttt{\&}（与）、\texttt{|}（或）、\texttt{\textasciicircum}（异或）
			      \item 逻辑非：!（一元运算符）
			      \item 优先级：! > \&\& > ||
		      \end{itemize}
	      \end{bluebox}

	      \vspace{1em}

	\item \textbf{以下代码的输出结果是}
	      \begin{lstlisting}[language=Java]
boolean flag = false;
if (flag) {
    System.out.println("True");
} else {
    System.out.println("False");
}
	      \end{lstlisting}
	      \begin{enumerate}[label=\Alph*.]
		      \item True
		      \item False
		      \item 编译错误
		      \item 运行时错误
	      \end{enumerate}

	      \begin{bluebox}{答案与深度解析}
		      \textbf{答案：B. False}

		      \textbf{【核心认知框架>}
		      \begin{itemize}[label={}, leftmargin=*, nosep]
			      \item \textbf{题目本质}：boolean类型与条件语句执行流程
		      \end{itemize}

		      \textbf{【选项原理深度拆解】}

		      \begin{itemize}[leftmargin=*, nosep, label={}]
			      \item \textbf{代码分析}：
			      \begin{itemize}[label={}, nosep]
				      \item flag = false
				      \item if(false) 不执行if块
				      \item 执行else块，输出"False"
			      \end{itemize}

			      \item[A] \textbf{True — 错误}
			      \item[B] \textbf{False — 正确}
			      \item[C] \textbf{编译错误 — 错误（语法正确）}
			      \item[D] \textbf{运行时错误 — 错误（无异常）}
		      \end{itemize}

		      \textbf{【应背诵知识点>}
		      \begin{itemize}[label={}, nosep]
			      \item Java中boolean是独立类型，不是int
			      \item if条件必须是boolean类型
			      \item Java不支持"非零即真"
		      \end{itemize}
	      \end{bluebox}

	      \vspace{1em}

	\item \textbf{在Java中，for循环的三个表达式都可以省略吗？}
	      \begin{enumerate}[label=\Alph*.]
		      \item 可以
		      \item 不可以
		      \item 第一个和第三个可以省略，第二个不可以
		      \item 第一个和第二个可以省略，第三个不可以
	      \end{enumerate}

	      \begin{bluebox}{答案与深度解析}
		      \textbf{答案：A. 可以}

		      \textbf{【核心认知框架>}
		      \begin{itemize}[label={}, leftmargin=*, nosep]
			      \item \textbf{题目本质}：for循环可省略语法特性
		      \end{itemize}

		      \textbf{【选项原理深度拆解】}

		      \begin{itemize}[leftmargin=*, nosep, label={}]
			      \item[A] \textbf{可以 — 正确}
			      \begin{itemize}[label={}, nosep]
				      \item JLS §14.14.1：三个表达式均可省略
				      \item 分号不能省略
				      \item for(;;) 等价于 while(true)
			      \end{itemize}

			      \item[B] \textbf{不可以 — 错误}
			      \item[C/D] \textbf{部分省略 — 错误}
		      \end{itemize}

		      \textbf{【应背诵知识点>}
		      \begin{itemize}[label={}, nosep]
			      \item for(初始化; 条件; 更新) \{ \}
			      \item 全部省略：for(;;)
			      \item 死循环写法：for(;;) 或 while(true)
		      \end{itemize}
	      \end{bluebox}

	      \vspace{1em}

	\item \textbf{以下关于数组的说法错误的是}
	      \begin{enumerate}[label=\Alph*.]
		      \item 数组是相同类型元素的有序集合
		      \item 数组的长度在创建后不能改变
		      \item 数组可以存储不同类型的元素
		      \item 可以通过索引访问数组元素
	      \end{enumerate}

	      \begin{bluebox}{答案与深度解析}
		      \textbf{答案：C. 数组可以存储不同类型的元素}

		      \textbf{【核心认知框架>}
		      \begin{itemize}[label={}, leftmargin=*, nosep]
			      \item \textbf{题目本质}：数组核心特性——同类型、固定长度

		      \textbf{【选项原理深度拆解】}

				\item \textbf{A. 相同类型 — 正确}
				\begin{itemize}[label={}, nosep]
					\item 数组只能在创建时指定类型
					\item 只能存储该类型元素
				\end{itemize}

				\item \textbf{B. 长度固定 — 正确}
				\begin{itemize}[label={}, nosep]
					\item 数组长度在创建时确定
					\item 使用length属性获取长度
				\end{itemize}

				\item \textbf{C. 不同类型 — 错误（本题答案）}
				\begin{itemize}[label={}, nosep]
					\item 数组只能存储相同类型元素
					\item 若需存储不同类型，使用集合或Object[]
				\end{itemize}

				\item \textbf{D. 索引访问 — 正确}
				\begin{itemize}[label={}, nosep]
					\item 索引从0开始
					\item 最大索引 = length - 1
				\end{itemize}

				\textbf{【应背诵知识点>}
				\begin{itemize}[label={}, nosep]
					\item 数组创建：new int[5] 或 \{1,2,3\}
					\item 索引范围：0 ~ length-1
					\item 数组是对象，有length属性
				\end{itemize}
		      \end{itemize}

	      \end{bluebox}

	      \vspace{1em}

	\item \textbf{以下关于Java中的可变参数的说法错误的是}
	      \begin{enumerate}[label=\Alph*.]
		      \item 可变参数必须是方法的最后一个参数
		      \item 可变参数可以接受零个或多个参数
		      \item 可变参数在方法内部被视为数组
		      \item 可变参数可以和普通参数一起使用，且位置任意
	      \end{enumerate}

	      \begin{bluebox}{答案与深度解析}
		      \textbf{答案：D. 可变参数可以和普通参数一起使用，且位置任意}

		      \textbf{【核心认知框架>}
		      \begin{itemize}[label={}, leftmargin=*, nosep]
			      \item \textbf{题目本质}：可变参数使用规则

		      \textbf{【选项原理深度拆解】}

		      \item \textbf{A. 最后参数 — 正确}
		      \begin{itemize}[label={}, nosep]
			      \item 可变参数必须是最后一个参数
		      \end{itemize}

		      \item \textbf{B. 零或多 — 正确}
		      \begin{itemize}[label={}, nosep]
			      \item 可以不传参数，也可以传多个
		      \end{itemize}

		      \item \textbf{C. 视为数组 — 正确}
		      \begin{itemize}[label={}, nosep]
			      \item 方法内当作数组使用
			      \item 可以用for循环遍历
		      \end{itemize}

		      \item \textbf{D. 位置任意 — 错误（本题答案）}
		      \begin{itemize}[label={}, nosep]
			      \item 可变参数必须是最后一个参数
			      \item 不能有多个可变参数
		      \end{itemize}

		      \textbf{【应背诵知识点>}
		      \begin{itemize}[label={}, nosep]
			      \item 语法：方法名(类型... 参数名)
			      \item 只能是最后一个参数
			      \item 方法重载时可能有歧义
		      \end{itemize}
		      \end{itemize}
	      \end{bluebox}

\end{enumerate}

\end{document}
